\chapter{Unit 5: Multivariate Differential Calculus and Constrained Extrema}

\section{Section A: Multiple Choice Practice Questions (60 Curated Problems)}

\noindent\textbf{Q1. For a limit of a function of two variables f ( x , y ) to exist as ( x , y ) $\to$ ( a , b ) , the limit must be:} \hfill [Tag: limit, continuity and differentiability of functions of two variables | Easy]
\begin{enumerate}[label=(\alph*),topsep=2pt,itemsep=1pt,partopsep=0pt,parsep=0pt]
    \item {{Different along different paths}}
    \item {{Zero along all paths}}
    \item {{Infinite along at least one path}}
    \item {{The same along any path approaching ( a , b )}}
\end{enumerate}
\vspace{2mm}

\noindent\textbf{Q2. A function f ( x , y ) is said to be continuous at a point ( a , b ) if:} \hfill [Tag: limit, continuity and differentiability of functions of two variables | Easy]
\begin{enumerate}[label=(\alph*),topsep=2pt,itemsep=1pt,partopsep=0pt,parsep=0pt]
    \item {{lim ( x , y ) $\to$ ( a , b )}}
    \item {{f ( a , b )}}
    \item {{lim ( x , y ) $\to$ ( a , b )}}
    \item {{Its partial derivatives evaluate to zero}}
\end{enumerate}
\vspace{2mm}

\noindent\textbf{Q3. If a function f ( x , y ) is differentiable at a point ( a , b ) , then at that point it must also be:} \hfill [Tag: limit, continuity and differentiability of functions of two variables | Easy]
\begin{enumerate}[label=(\alph*),topsep=2pt,itemsep=1pt,partopsep=0pt,parsep=0pt]
    \item {{Undefined}}
    \item {{Continuous}}
    \item {{Equal to zero}}
    \item {{A constant function}}
\end{enumerate}
\vspace{2mm}

\noindent\textbf{Q4. If z = f ( x , y ) where x = u ( t ) and y = v ( t ) , what is the total derivative d t d z according to the chain rule?} \hfill [Tag: chain rule | Easy]
\begin{enumerate}[label=(\alph*),topsep=2pt,itemsep=1pt,partopsep=0pt,parsep=0pt]
    \item {{d t}}
    \item {{$\partial$ x}}
    \item {{$\partial$ x}}
    \item {{$\partial$ x}}
\end{enumerate}
\vspace{2mm}

\noindent\textbf{Q5. If u = f ( x , y ) and x = g ( r , s ) , y = h ( r , s ) , how is the partial derivative $\partial$ r $\partial$ u expressed?} \hfill [Tag: chain rule | Easy]
\begin{enumerate}[label=(\alph*),topsep=2pt,itemsep=1pt,partopsep=0pt,parsep=0pt]
    \item {{$\partial$ x}}
    \item {{$\partial$ x}}
    \item {{$\partial$ r}}
    \item {{$\partial$ x}}
\end{enumerate}
\vspace{2mm}

\noindent\textbf{Q6. When changing variables in a double integral from ( x , y ) to ( u , v ) , the area element d x d y becomes:} \hfill [Tag: change of variables | Easy]
\begin{enumerate}[label=(\alph*),topsep=2pt,itemsep=1pt,partopsep=0pt,parsep=0pt]
    \item {{| J |}}
    \item {{J 2}}
    \item {{| J | d u d v}}
    \item {{d u d v}}
\end{enumerate}
\vspace{2mm}

\noindent\textbf{Q7. In the standard change of variables to polar coordinates, what are x and y substituted with?} \hfill [Tag: change of variables | Easy]
\begin{enumerate}[label=(\alph*),topsep=2pt,itemsep=1pt,partopsep=0pt,parsep=0pt]
    \item {{x =}}
    \item {{x =}}
    \item {{x =}}
    \item {{x =}}
\end{enumerate}
\vspace{2mm}

\noindent\textbf{Q8. A function f ( x , y ) is homogeneous of degree n if for any scalar t > 0 :} \hfill [Tag: Euler's theorem for homogeneous equations | Easy]
\begin{enumerate}[label=(\alph*),topsep=2pt,itemsep=1pt,partopsep=0pt,parsep=0pt]
    \item {{f ( t x , t y ) =}}
    \item {{f ( t x , t y ) =}}
    \item {{f ( t x , t y ) =}}
    \item {{f ( t x , t y ) =}}
\end{enumerate}
\vspace{2mm}

\noindent\textbf{Q9. According to Euler's theorem, if u ( x , y ) is a homogeneous function of degree n , then x $\partial$ x $\partial$ u + y $\partial$ y $\partial$ u is equal to:} \hfill [Tag: Euler's theorem for homogeneous equations | Easy]
\begin{enumerate}[label=(\alph*),topsep=2pt,itemsep=1pt,partopsep=0pt,parsep=0pt]
    \item {{n u 2}}
    \item {{n u}}
    \item {{u n}}
    \item {{0}}
\end{enumerate}
\vspace{2mm}

\noindent\textbf{Q10. What is the degree of homogeneity of the function f ( x , y ) = x + y x 3 + y 3 ?} \hfill [Tag: Euler's theorem for homogeneous equations | Easy]
\begin{enumerate}[label=(\alph*),topsep=2pt,itemsep=1pt,partopsep=0pt,parsep=0pt]
    \item {{1}}
    \item {{2}}
    \item {{3}}
    \item {{4}}
\end{enumerate}
\vspace{2mm}

\noindent\textbf{Q11. The Jacobian of u and v with respect to x and y , denoted by J = $\partial$ ( x , y ) $\partial$ ( u , v ) , is defined as:} \hfill [Tag: Jacobians | Easy]
\begin{enumerate}[label=(\alph*),topsep=2pt,itemsep=1pt,partopsep=0pt,parsep=0pt]
    \item {{The matrix (}}
    \item {{The sum of the partial derivatives u x}}
    \item {{The product of u x}}
    \item {{The determinant of the matrix (}}
\end{enumerate}
\vspace{2mm}

\noindent\textbf{Q12. If J is the Jacobian of ( u , v ) with respect to ( x , y ) , and J ' is the Jacobian of ( x , y ) with respect to ( u , v ) , then J $\times$ J ' equals:} \hfill [Tag: Jacobians | Easy]
\begin{enumerate}[label=(\alph*),topsep=2pt,itemsep=1pt,partopsep=0pt,parsep=0pt]
    \item {{uvx y}}
    \item {{1}}
    \item {{-1}}
    \item {{0}}
\end{enumerate}
\vspace{2mm}

\noindent\textbf{Q13. What is the Jacobian $\partial$ ( r , $\theta$ ) $\partial$ ( x , y ) for the polar coordinate transformation x = r cos $\theta$ , y = r sin $\theta$ ?} \hfill [Tag: Jacobians | Easy]
\begin{enumerate}[label=(\alph*),topsep=2pt,itemsep=1pt,partopsep=0pt,parsep=0pt]
    \item {{r 2}}
    \item {{r}}
    \item {{sin $\theta$ cos $\theta$}}
    \item {{1}}
\end{enumerate}
\vspace{2mm}

\noindent\textbf{Q14. A point ( a , b ) is a critical (stationary) point of a differentiable function f ( x , y ) if:} \hfill [Tag: extrema of functions of two variables | Easy]
\begin{enumerate}[label=(\alph*),topsep=2pt,itemsep=1pt,partopsep=0pt,parsep=0pt]
    \item {{f ( a , b ) =}}
    \item {{f x}}
    \item {{f xx}}
    \item {{f x}}
\end{enumerate}
\vspace{2mm}

\noindent\textbf{Q15. Let r = f xx , s = f x y , and t = f yy . For a critical point to be a local minimum, which conditions must hold?} \hfill [Tag: extrema of functions of two variables | Easy]
\begin{enumerate}[label=(\alph*),topsep=2pt,itemsep=1pt,partopsep=0pt,parsep=0pt]
    \item {{r t -}}
    \item {{r t -}}
    \item {{r t -}}
    \item {{r t -}}
\end{enumerate}
\vspace{2mm}

\noindent\textbf{Q16. If r t - s 2 < 0 at a critical point of f ( x , y ) , the point is called a:} \hfill [Tag: extrema of functions of two variables | Easy]
\begin{enumerate}[label=(\alph*),topsep=2pt,itemsep=1pt,partopsep=0pt,parsep=0pt]
    \item {{Saddle point}}
    \item {{Local maximum}}
    \item {{Local minimum}}
    \item {{Point of inflection}}
\end{enumerate}
\vspace{2mm}

\noindent\textbf{Q17. If r t - s 2 = 0 at a critical point, what does the second derivative test conclude?} \hfill [Tag: extrema of functions of two variables | Easy]
\begin{enumerate}[label=(\alph*),topsep=2pt,itemsep=1pt,partopsep=0pt,parsep=0pt]
    \item {{The test is inconclusive}}
    \item {{It is definitively a saddle point}}
    \item {{It is definitively a local minimum}}
    \item {{It is definitively a local maximum}}
\end{enumerate}
\vspace{2mm}

\noindent\textbf{Q18. Lagrange's method of undetermined multipliers is primarily used for:} \hfill [Tag: Lagrange's method of undetermined multipliers | Easy]
\begin{enumerate}[label=(\alph*),topsep=2pt,itemsep=1pt,partopsep=0pt,parsep=0pt]
    \item {{Finding the limit of a multivariable function}}
    \item {{Finding extrema of a function subject to equality constraints}}
    \item {{Calculating the Jacobian of a transformation}}
    \item {{Testing for continuity}}
\end{enumerate}
\vspace{2mm}

\noindent\textbf{Q19. To find the extrema of f ( x , y ) subject to the constraint g ( x , y ) = 0 , we form the Lagrangian function F ( x , y , $\lambda$ ) as:} \hfill [Tag: Lagrange's method of undetermined multipliers | Easy]
\begin{enumerate}[label=(\alph*),topsep=2pt,itemsep=1pt,partopsep=0pt,parsep=0pt]
    \item {{F ( x , y , $\lambda$ ) =}}
    \item {{F ( x , y , $\lambda$ ) =}}
    \item {{F ( x , y , $\lambda$ ) =}}
    \item {{F ( x , y , $\lambda$ ) =}}
\end{enumerate}
\vspace{2mm}

\noindent\textbf{Q20. In the Lagrangian function F ( x , y , $\lambda$ ) = f ( x , y ) + $\lambda$ g ( x , y ) , what is the variable $\lambda$ called?} \hfill [Tag: Lagrange's method of undetermined multipliers | Easy]
\begin{enumerate}[label=(\alph*),topsep=2pt,itemsep=1pt,partopsep=0pt,parsep=0pt]
    \item {{Euler constant}}
    \item {{Jacobian determinant}}
    \item {{Critical variable}}
    \item {{Lagrange multiplier}}
\end{enumerate}
\vspace{2mm}

\noindent\textbf{Q21. Evaluate the limit: lim ( x , y ) $\to$ ( 0 , 0 ) x 2 + y 2 x 2 y .} \hfill [Tag: limit, continuity and differentiability of functions of two variables | Medium]
\begin{enumerate}[label=(\alph*),topsep=2pt,itemsep=1pt,partopsep=0pt,parsep=0pt]
    \item {{1}}
    \item {{2}}
    \item {{0}}
    \item {{Does not exist}}
\end{enumerate}
\vspace{2mm}

\noindent\textbf{Q22. Determine the limit lim ( x , y ) $\to$ ( 0 , 0 ) x 2 + y 2 x y along the path y = m x .} \hfill [Tag: limit, continuity and differentiability of functions of two variables | Medium]
\begin{enumerate}[label=(\alph*),topsep=2pt,itemsep=1pt,partopsep=0pt,parsep=0pt]
    \item {{Does not exist}}
    \item {{0}}
    \item {{1 + m 2}}
    \item {{m}}
\end{enumerate}
\vspace{2mm}

\noindent\textbf{Q23. If z = x 2 + y 2 where x = e t cos t and y = e t sin t , find d t d z .} \hfill [Tag: chain rule | Medium]
\begin{enumerate}[label=(\alph*),topsep=2pt,itemsep=1pt,partopsep=0pt,parsep=0pt]
    \item {{2 e 2 t}}
    \item {{e 2 t}}
    \item {{2 e t}}
    \item {{4 e 2 t}}
\end{enumerate}
\vspace{2mm}

\noindent\textbf{Q24. Let u = f ( x - y , y - z , z - x ) . Find $\partial$ x $\partial$ u + $\partial$ y $\partial$ u + $\partial$ z $\partial$ u .} \hfill [Tag: chain rule | Medium]
\begin{enumerate}[label=(\alph*),topsep=2pt,itemsep=1pt,partopsep=0pt,parsep=0pt]
    \item {{0}}
    \item {{- 1}}
    \item {{1}}
    \item {{3}}
\end{enumerate}
\vspace{2mm}

\noindent\textbf{Q25. If u = sin - 1 ( x + y x 2 + y 2 ) , find the value of x $\partial$ x $\partial$ u + y $\partial$ y $\partial$ u .} \hfill [Tag: Euler's theorem for homogeneous equations | Medium]
\begin{enumerate}[label=(\alph*),topsep=2pt,itemsep=1pt,partopsep=0pt,parsep=0pt]
    \item {{cot u}}
    \item {{tan u}}
    \item {{sin u}}
    \item {{cos u}}
\end{enumerate}
\vspace{2mm}

\noindent\textbf{Q26. Given u = lo g ( x + y x 4 + y 4 ) , evaluate x $\partial$ x $\partial$ u + y $\partial$ y $\partial$ u .} \hfill [Tag: Euler's theorem for homogeneous equations | Medium]
\begin{enumerate}[label=(\alph*),topsep=2pt,itemsep=1pt,partopsep=0pt,parsep=0pt]
    \item {{e 3}}
    \item {{3}}
    \item {{4}}
    \item {{lo g}}
\end{enumerate}
\vspace{2mm}

\noindent\textbf{Q27. If x = r cos $\theta$ and y = r sin $\theta$ , find the Jacobian J = $\partial$ ( r , $\theta$ ) $\partial$ ( x , y ) .} \hfill [Tag: Jacobians | Medium]
\begin{enumerate}[label=(\alph*),topsep=2pt,itemsep=1pt,partopsep=0pt,parsep=0pt]
    \item {{1}}
    \item {{r 2}}
    \item {{sin $\theta$ cos $\theta$}}
    \item {{r}}
\end{enumerate}
\vspace{2mm}

\noindent\textbf{Q28. Given u = x + y and v = x y , find the Jacobian $\partial$ ( x , y ) $\partial$ ( u , v ) .} \hfill [Tag: Jacobians | Medium]
\begin{enumerate}[label=(\alph*),topsep=2pt,itemsep=1pt,partopsep=0pt,parsep=0pt]
    \item {{x -}}
    \item {{x +}}
    \item {{1}}
    \item {{y -}}
\end{enumerate}
\vspace{2mm}

\noindent\textbf{Q29. Under the transformation u = x - y , v = x + y , what does the differential operator $\partial$ x 2 $\partial$ 2 - $\partial$ y 2 $\partial$ 2 become?} \hfill [Tag: change of variables | Medium]
\begin{enumerate}[label=(\alph*),topsep=2pt,itemsep=1pt,partopsep=0pt,parsep=0pt]
    \item {{4 $\partial$ u $\partial$ v}}
    \item {{2 $\partial$ u $\partial$ v}}
    \item {{$\partial$ u 2}}
    \item {{$\partial$ u $\partial$ v}}
\end{enumerate}
\vspace{2mm}

\noindent\textbf{Q30. Find the critical points of the function f ( x , y ) = x 2 + y 2 - 2 x - 4 y + 5 .} \hfill [Tag: extrema of functions of two variables | Medium]
\begin{enumerate}[label=(\alph*),topsep=2pt,itemsep=1pt,partopsep=0pt,parsep=0pt]
    \item {{( 1 , 2 )}}
    \item {{( 1 , - 2 )}}
    \item {{( 2 , 1 )}}
    \item {{( - 1 , 2 )}}
\end{enumerate}
\vspace{2mm}

\noindent\textbf{Q31. Determine the nature of the critical point ( 0 , 0 ) for the function f ( x , y ) = x 2 - y 2 .} \hfill [Tag: extrema of functions of two variables | Medium]
\begin{enumerate}[label=(\alph*),topsep=2pt,itemsep=1pt,partopsep=0pt,parsep=0pt]
    \item {{Inconclusive}}
    \item {{Local maximum}}
    \item {{Saddle point}}
    \item {{Local minimum}}
\end{enumerate}
\vspace{2mm}

\noindent\textbf{Q32. For a function f ( x , y ) , at a critical point ( a , b ) , let r = f xx , s = f x y , and t = f yy . Which condition guarantees a local maximum?} \hfill [Tag: extrema of functions of two variables | Medium]
\begin{enumerate}[label=(\alph*),topsep=2pt,itemsep=1pt,partopsep=0pt,parsep=0pt]
    \item {{r t -}}
    \item {{r t -}}
    \item {{r t -}}
    \item {{r t -}}
\end{enumerate}
\vspace{2mm}

\noindent\textbf{Q33. To find the extreme values of f ( x , y ) = x + y subject to x 2 + y 2 = 1 , what are the Lagrange multiplier equations?} \hfill [Tag: Lagrange's method of undetermined multipliers | Medium]
\begin{enumerate}[label=(\alph*),topsep=2pt,itemsep=1pt,partopsep=0pt,parsep=0pt]
    \item {{1 =}}
    \item {{2 x =}}
    \item {{x =}}
    \item {{1 =}}
\end{enumerate}
\vspace{2mm}

\noindent\textbf{Q34. Find the maximum value of f ( x , y ) = x y subject to the constraint x + y = 10 using Lagrange multipliers.} \hfill [Tag: Lagrange's method of undetermined multipliers | Medium]
\begin{enumerate}[label=(\alph*),topsep=2pt,itemsep=1pt,partopsep=0pt,parsep=0pt]
    \item {{50}}
    \item {{20}}
    \item {{25}}
    \item {{100}}
\end{enumerate}
\vspace{2mm}

\noindent\textbf{Q35. Which of the following functions is continuous at ( 0 , 0 ) ?} \hfill [Tag: limit, continuity and differentiability of functions of two variables | Medium]
\begin{enumerate}[label=(\alph*),topsep=2pt,itemsep=1pt,partopsep=0pt,parsep=0pt]
    \item {{f ( x , y ) =}}
    \item {{f ( x , y ) =}}
    \item {{f ( x , y ) =}}
    \item {{f ( x , y ) =}}
\end{enumerate}
\vspace{2mm}

\noindent\textbf{Q36. If w = f ( x , y ) , where x = u 2 - v 2 and y = 2 uv , express $\partial$ u $\partial$ w in terms of w x and w y .} \hfill [Tag: chain rule | Medium]
\begin{enumerate}[label=(\alph*),topsep=2pt,itemsep=1pt,partopsep=0pt,parsep=0pt]
    \item {{2 v w x}}
    \item {{- 2 v w x}}
    \item {{2 u w x}}
    \item {{2 u w x}}
\end{enumerate}
\vspace{2mm}

\noindent\textbf{Q37. If u ( x , y ) is a homogeneous function of degree n , what is the degree of homogeneity of $\partial$ x $\partial$ u ?} \hfill [Tag: Euler's theorem for homogeneous equations | Medium]
\begin{enumerate}[label=(\alph*),topsep=2pt,itemsep=1pt,partopsep=0pt,parsep=0pt]
    \item {{1}}
    \item {{n -}}
    \item {{n +}}
    \item {{n}}
\end{enumerate}
\vspace{2mm}

\noindent\textbf{Q38. Let u = y x and v = z y , w = x z . What is the value of the Jacobian $\partial$ ( x , y , z ) $\partial$ ( u , v , w ) ?} \hfill [Tag: Jacobians | Medium]
\begin{enumerate}[label=(\alph*),topsep=2pt,itemsep=1pt,partopsep=0pt,parsep=0pt]
    \item {{x yz}}
    \item {{1}}
    \item {{- 1}}
    \item {{0}}
\end{enumerate}
\vspace{2mm}

\noindent\textbf{Q39. For f ( x , y ) = x 3 + y 3 - 3 a x y , finding critical points involves solving:} \hfill [Tag: extrema of functions of two variables | Medium]
\begin{enumerate}[label=(\alph*),topsep=2pt,itemsep=1pt,partopsep=0pt,parsep=0pt]
    \item {{3 x 2}}
    \item {{x 2}}
    \item {{x 3}}
    \item {{3 x 2}}
\end{enumerate}
\vspace{2mm}

\noindent\textbf{Q40. When applying Lagrange multipliers to minimize f ( x , y ) subject to two constraints g ( x , y ) = 0 and h ( x , y ) = 0 , how many multipliers are needed?} \hfill [Tag: Lagrange's method of undetermined multipliers | Medium]
\begin{enumerate}[label=(\alph*),topsep=2pt,itemsep=1pt,partopsep=0pt,parsep=0pt]
    \item {{One}}
    \item {{None}}
    \item {{Two}}
    \item {{Three}}
\end{enumerate}
\vspace{2mm}

\noindent\textbf{Q41. Consider the function f ( x , y ) = x 6 + y 2 x 3 y for ( x , y )  = ( 0 , 0 ) and f ( 0 , 0 ) = 0 . Which of the following statements is true regarding its differentiability at ( 0 , 0 ) ?} \hfill [Tag: limit, continuity and differentiability of functions of two variables | Hard]
\begin{enumerate}[label=(\alph*),topsep=2pt,itemsep=1pt,partopsep=0pt,parsep=0pt]
    \item {{f}}
    \item {{The directional derivatives do not exist in all directions at ( 0 , 0 )}}
    \item {{f}}
    \item {{The directional derivatives exist in all directions, but f}}
\end{enumerate}
\vspace{2mm}

\noindent\textbf{Q42. Let f ( x , y ) = ( x 2 + y 2 ) sin ( x 2 + y 2 1 ) for ( x , y )  = ( 0 , 0 ) , and 0 otherwise. Which of the following holds at ( 0 , 0 ) ?} \hfill [Tag: limit, continuity and differentiability of functions of two variables | Hard]
\begin{enumerate}[label=(\alph*),topsep=2pt,itemsep=1pt,partopsep=0pt,parsep=0pt]
    \item {{f}}
    \item {{f}}
    \item {{f}}
    \item {{f}}
\end{enumerate}
\vspace{2mm}

\noindent\textbf{Q43. Let z = f ( x , y ) , where x = u 2 - v 2 and y = 2 uv . If f satisfies Laplace's equation $\partial$ x 2 $\partial$ 2 z + $\partial$ y 2 $\partial$ 2 z = 0 , what does $\partial$ u 2 $\partial$ 2 z + $\partial$ v 2 $\partial$ 2 z equal?} \hfill [Tag: chain rule | Hard]
\begin{enumerate}[label=(\alph*),topsep=2pt,itemsep=1pt,partopsep=0pt,parsep=0pt]
    \item {{4 ( u 2}}
    \item {{( u 2}}
    \item {{0}}
    \item {{4 ( u 2}}
\end{enumerate}
\vspace{2mm}

\noindent\textbf{Q44. If u = f ( x y y - x , x z z - x ) , what is the value of x 2 $\partial$ x $\partial$ u + y 2 $\partial$ y $\partial$ u + z 2 $\partial$ z $\partial$ u ?} \hfill [Tag: chain rule | Hard]
\begin{enumerate}[label=(\alph*),topsep=2pt,itemsep=1pt,partopsep=0pt,parsep=0pt]
    \item {{1}}
    \item {{x yz}}
    \item {{0}}
    \item {{x +}}
\end{enumerate}
\vspace{2mm}

\noindent\textbf{Q45. If u = sin - 1 ( x + y x 3 + y 3 ) , what is the value of x 2 $\partial$ x 2 $\partial$ 2 u + 2 x y $\partial$ x $\partial$ y $\partial$ 2 u + y 2 $\partial$ y 2 $\partial$ 2 u ?} \hfill [Tag: Euler's theorem for homogeneous equations | Hard]
\begin{enumerate}[label=(\alph*),topsep=2pt,itemsep=1pt,partopsep=0pt,parsep=0pt]
    \item {{2 tan u ( 2 sec 2}}
    \item {{2 tan u}}
    \item {{2 tan u ( 2 tan 2}}
    \item {{2 tan u ( 2 sec 2}}
\end{enumerate}
\vspace{2mm}

\noindent\textbf{Q46. Let u = lo g ( x + y x 4 + y 4 ) . Evaluate x $\partial$ x $\partial$ u + y $\partial$ y $\partial$ u .} \hfill [Tag: Euler's theorem for homogeneous equations | Hard]
\begin{enumerate}[label=(\alph*),topsep=2pt,itemsep=1pt,partopsep=0pt,parsep=0pt]
    \item {{e 3}}
    \item {{3 u}}
    \item {{3}}
    \item {{u}}
\end{enumerate}
\vspace{2mm}

\noindent\textbf{Q47. If x + y + z = u , y + z = uv , and z = uv w , find the Jacobian $\partial$ ( u , v , w ) $\partial$ ( x , y , z ) .} \hfill [Tag: Jacobians | Hard]
\begin{enumerate}[label=(\alph*),topsep=2pt,itemsep=1pt,partopsep=0pt,parsep=0pt]
    \item {{u 2}}
    \item {{u 2}}
    \item {{u v 2}}
    \item {{u 2}}
\end{enumerate}
\vspace{2mm}

\noindent\textbf{Q48. Given functions u = x + y + z , v = x 2 + y 2 + z 2 , and w = x 3 + y 3 + z 3 - 3 x yz , what is the relationship between u , v , w ?} \hfill [Tag: Jacobians | Hard]
\begin{enumerate}[label=(\alph*),topsep=2pt,itemsep=1pt,partopsep=0pt,parsep=0pt]
    \item {{They are functionally independent.}}
    \item {{w =}}
    \item {{w =}}
    \item {{w =}}
\end{enumerate}
\vspace{2mm}

\noindent\textbf{Q49. Find the nature of the critical point ( 0 , 0 ) for the function f ( x , y ) = x 4 + y 4 - 2 x 2 + 4 x y - 2 y 2 .} \hfill [Tag: extrema of functions of two variables | Hard]
\begin{enumerate}[label=(\alph*),topsep=2pt,itemsep=1pt,partopsep=0pt,parsep=0pt]
    \item {{Inconclusive by the second derivative test, but it is a saddle point}}
    \item {{Local minimum}}
    \item {{Local maximum}}
    \item {{Saddle point}}
\end{enumerate}
\vspace{2mm}

\noindent\textbf{Q50. What are the absolute extrema of f ( x , y ) = x 2 y on the closed domain D = {( x , y ) | x 2 + y 2 $\le$ 1 } ?} \hfill [Tag: extrema of functions of two variables | Hard]
\begin{enumerate}[label=(\alph*),topsep=2pt,itemsep=1pt,partopsep=0pt,parsep=0pt]
    \item {{Maximum is 2}}
    \item {{Maximum is 3 3}}
    \item {{Maximum is 9}}
    \item {{Maximum is 9}}
\end{enumerate}
\vspace{2mm}

\noindent\textbf{Q51. Using Lagrange multipliers, find the minimum distance from the origin to the surface x 2 y 2 z 2 = a 6 , where a > 0 .} \hfill [Tag: Lagrange's method of undetermined multipliers | Hard]
\begin{enumerate}[label=(\alph*),topsep=2pt,itemsep=1pt,partopsep=0pt,parsep=0pt]
    \item {{a 3}}
    \item {{3}}
    \item {{3 a}}
    \item {{3}}
\end{enumerate}
\vspace{2mm}

\noindent\textbf{Q52. Find the maximum value of f ( x , y , z ) = x + y + z subject to two constraints: x 2 + y 2 = 2 and x + z = 1 .} \hfill [Tag: Lagrange's method of undetermined multipliers | Hard]
\begin{enumerate}[label=(\alph*),topsep=2pt,itemsep=1pt,partopsep=0pt,parsep=0pt]
    \item {{1 +}}
    \item {{3}}
    \item {{1 +}}
    \item {{1 +}}
\end{enumerate}
\vspace{2mm}

\noindent\textbf{Q53. Under the transformation x = u cos v , y = u sin v , the Laplacian operator $\Delta$ f = $\partial$ x 2 $\partial$ 2 f + $\partial$ y 2 $\partial$ 2 f transforms to:} \hfill [Tag: change of variables | Hard]
\begin{enumerate}[label=(\alph*),topsep=2pt,itemsep=1pt,partopsep=0pt,parsep=0pt]
    \item {{$\partial$ u 2}}
    \item {{$\partial$ u 2}}
    \item {{$\partial$ u 2}}
    \item {{$\partial$ u 2}}
\end{enumerate}
\vspace{2mm}

\noindent\textbf{Q54. Let u = x - y and v = x + y . Transform the differential equation $\partial$ x 2 $\partial$ 2 z - $\partial$ y 2 $\partial$ 2 z = 0 into the ( u , v ) coordinate system.} \hfill [Tag: change of variables | Hard]
\begin{enumerate}[label=(\alph*),topsep=2pt,itemsep=1pt,partopsep=0pt,parsep=0pt]
    \item {{$\partial$ u $\partial$ v}}
    \item {{$\partial$ u $\partial$ v}}
    \item {{$\partial$ u 2}}
    \item {{$\partial$ u 2}}
\end{enumerate}
\vspace{2mm}

\noindent\textbf{Q55. Let f ( x , y ) = x 2 + y 2 x y ( x 2 - y 2 ) for ( x , y )  = ( 0 , 0 ) and 0 at ( 0 , 0 ) . What is the relationship between the mixed partial derivatives at ( 0 , 0 ) ?} \hfill [Tag: limit, continuity and differentiability of functions of two variables | Hard]
\begin{enumerate}[label=(\alph*),topsep=2pt,itemsep=1pt,partopsep=0pt,parsep=0pt]
    \item {{f x y}}
    \item {{f x y}}
    \item {{f x y}}
    \item {{The mixed partial derivatives do not exist at ( 0 , 0 )}}
\end{enumerate}
\vspace{2mm}

\noindent\textbf{Q56. Suppose F ( x , y , z ) = 0 defines z implicitly as a function of x and y . What is the expression for $\partial$ x 2 $\partial$ 2 z ?} \hfill [Tag: chain rule | Hard]
\begin{enumerate}[label=(\alph*),topsep=2pt,itemsep=1pt,partopsep=0pt,parsep=0pt]
    \item {{F z}}
    \item {{- F z}}
    \item {{- F z}}
    \item {{- F z}}
\end{enumerate}
\vspace{2mm}

\noindent\textbf{Q57. If u ( x , y ) is a homogeneous function of degree n , what is the value of x 3 u xxx + 3 x 2 y u xx y + 3 x y 2 u x yy + y 3 u yyy ?} \hfill [Tag: Euler's theorem for homogeneous equations | Hard]
\begin{enumerate}[label=(\alph*),topsep=2pt,itemsep=1pt,partopsep=0pt,parsep=0pt]
    \item {{n 3}}
    \item {{n ( n -}}
    \item {{n ( n +}}
    \item {{n ( n -}}
\end{enumerate}
\vspace{2mm}

\noindent\textbf{Q58. Let u , v , w be the roots of the cubic equation $\lambda$ 3 - x $\lambda$ 2 + y $\lambda$ - z = 0 . What is the absolute value of the Jacobian | $\partial$ ( x , y , z ) $\partial$ ( u , v , w ) | ?} \hfill [Tag: Jacobians | Hard]
\begin{enumerate}[label=(\alph*),topsep=2pt,itemsep=1pt,partopsep=0pt,parsep=0pt]
    \item {{| ( u - v ) ( v - w ) ( w - u ) |}}
    \item {{| ( u -}}
    \item {{| ( u + v ) ( v + w ) ( w + u ) |}}
    \item {{1}}
\end{enumerate}
\vspace{2mm}

\noindent\textbf{Q59. Find the maximum volume of a rectangular parallelepiped inscribed in the ellipsoid a 2 x 2 + b 2 y 2 + c 2 z 2 = 1 with sides parallel to the coordinate axes.} \hfill [Tag: extrema of functions of two variables | Hard]
\begin{enumerate}[label=(\alph*),topsep=2pt,itemsep=1pt,partopsep=0pt,parsep=0pt]
    \item {{3}}
    \item {{3 3}}
    \item {{27}}
    \item {{3 3}}
\end{enumerate}
\vspace{2mm}

\noindent\textbf{Q60. Given a triangle with constant perimeter 2 s . What are the lengths of the sides a , b , c that maximize its area?} \hfill [Tag: Lagrange's method of undetermined multipliers | Hard]
\begin{enumerate}[label=(\alph*),topsep=2pt,itemsep=1pt,partopsep=0pt,parsep=0pt]
    \item {{a =}}
    \item {{a =}}
    \item {{a =}}
    \item {{The maximum area cannot be uniquely determined by the perimeter alone.}}
\end{enumerate}
\vspace{2mm}


\section{Section B: Comprehensive Subjective Problem with Detailed Analytical Rubric}

\begin{subjectiveq}{5}{1}
A rectangular box open at the top is to have a volume of 32 cubic units. Find the dimensions of the box that require the least material for its construction using Lagrange multipliers.
\end{subjectiveq}

\begin{solbox}
Let the dimensions be $x, y$ (base) and $z$ (height). Surface area $S(x,y,z) = xy + 2xz + 2yz$ subject to constraint $g(x,y,z) = xyz - 32 = 0$. Lagrangian: $F = xy + 2xz + 2yz + \lambda(xyz - 32)$. Setting partial derivatives to zero: $F_x = y + 2z + \lambda yz = 0$, $F_y = x + 2z + \lambda xz = 0$, $F_z = 2x + 2y + \lambda xy = 0$. From symmetry, $x = y = 2z$. Substituting into $xyz = 32 \implies (2z)(2z)(z) = 4z^3 = 32 \implies z = 2$. Thus $x = 4$, $y = 4$, $z = 2$ units.
\end{solbox}

\begin{examinertip}
Step 1: Identify the foundational definition and state governing theorems.\\
Step 2: Clearly write each algebraic transformation and boundary substitution.\\
Step 3: State the final numerical/algebraic invariant clearly in a concluding box.
\end{examinertip}
